
\begin{frame}{Appendix --- Campaign-Spending Attrition}
\hypertarget{app:spending}{}
\hypertarget{src:mechanism}{}
\small
The resource-attrition account holds that litigation drains a challenger's campaign funds faster than the front-runner's, tilting the contest before a vote is cast. We test it with a 2SLS of judicialization on the 2020$\to$2024 change in TSE campaign spending (SPCE contracted expenses), cut by vote rank (front-runner versus challenger) and by seat type. Attrition predicts that the challenger's spending, and their share of the top two, falls under exposure, most sharply where an incumbent is defended. Baseline controls; state FE; state-clustered; first stage $F=\FinFCont$ in contested seats, \FinFOpen{} in open seats and \FinFFull{} on all seats.
\smallskip
\begin{center}
\scriptsize
\resizebox{\linewidth}{!}{% Auto-generated by code/03_estimation/10_mechanism_finance.R
% Do not edit manually -- rerun the generating script to update

\begin{tabular}{l*{3}{c}}
\toprule\toprule
Subsample ($N$) & \textbf{$\Delta$ Log challenger spend} & \textbf{$\Delta$ Log front-runner spend} & \textbf{$\Delta$ Challenger spend share} \\
\midrule
\textbf{All seats} (5,560) & 0.086 & 0.102 & -0.027 \\
 & \textcolor{mygray}{(0.320)} & \textcolor{mygray}{(0.114)} & \textcolor{mygray}{(0.032)} \\
\textbf{Open seat} (2,026) & 0.285 & 0.100 & -0.023 \\
 & \textcolor{mygray}{(0.381)} & \textcolor{mygray}{(0.129)} & \textcolor{mygray}{(0.047)} \\
\rowcolor{mylight}
\textbf{Contested} (3,534) & -0.136 & 0.100 & -0.035 \\
\rowcolor{mylight}
 & \textcolor{mygray}{(0.379)} & \textcolor{mygray}{(0.141)} & \textcolor{mygray}{(0.028)} \\
\bottomrule\bottomrule
\end{tabular}
}
\end{center}
\smallskip
\footnotesize The evidence for attrition is suggestive at most. On the full sample every spending outcome is null, and in contested seats the signs line up (the challenger's spending falls, the front-runner's holds, the challenger's top-two share slips) but stay far from significant (best $p=\FinBestContP$). The channel we can measure is therefore voter withdrawal in these same contested seats \hyperlink{src:voterballot}{\beamergotobutton{ballot response}} rather than campaign finance: soft challenger support blanks out instead of being out-financed. \hfill\hyperlink{app:seatdisengage}{\beamerreturnbutton{Back}}
\end{frame}
